\documentclass{article}
\usepackage[utf8]{inputenc}

\title{Computación Distribuida \\ Tarea 04}
\author{Olea Olvera Marco Iván \\ Arteaga Hernández Alejandro Iabin}
\date{}

\usepackage[top=1.5cm, bottom=1.5cm, left=2cm, right=2cm]{geometry}
\usepackage{fouriernc}
\usepackage{amssymb}
\usepackage{listings}
\lstset{
  basicstyle=\ttfamily,
  mathescape
}

\usepackage{graphicx}
\graphicspath{{./}}

\begin{document}

\maketitle

\section{}

\section{}

Sea $\Pi = \{p_1,\ p_2,\ ...,\ p_n\}$ el conjunto de procesos. Supongamos que existe un protocolo que resuelve el problema (en particular, que se cumple \textit{acuerdo}). Este se diseña de tal manera que $p_1$ le manda un mensaje a los demás procesos avisando que decidirá 0. Pero como los canales de comunicación no son confiables, $p_1$ no tiene manera de saber si a los demás les llegó ese mensaje.

Si el protocolo se diseña de tal manera que $p_i$, $i \not= 1$, le tiene que confirmar a $p_1$ que recibió su mensaje, entonces $p_i$ no tiene manera de saber que sí le llegará su mensaje de confirmación a $p_1$. Evidentemente, sin importar cuántas rondas de confirmación se realicen, es imposible que un proceso sepa con certeza que el último mensaje que envío sí fue recibido.

Si existiera un protocolo que cumple \textit{acuerdo}, este se podría acortar eliminando el último mensaje enviado. Repitiendo esto iterativamente llegaríamos a un protocolo que no envía mensajes, es decir, un protocolo que cumple \textit{acuerdo} porque todos los procesos deciden un valor predeterminado, lo cual contradice el planteamiento del problema.

\section{}

En un sistema síncrono, cada proceso envía mensajes en un tiempo $t$ y estos mensajes son entregados en el tiempo $t+1$; a su vez, en el tiempo $t+1$ se envían mensajes que son entregados en el tiempo $t+2$ y así sucesivamente. En un sistema asíncrono los mensajes son entregados después de un retraso, el cual podría no estar acotado.

Otra diferencia es que en los sistemas asíncronos no se garantiza que los mensajes se reciben en el mismo orden en que se envían.

\section{}

\section{}

Las fuerzas de atracción gravitatoria que ejercen el Sol y la Luna sobre la Tierra causan la marea. \\ \\

\end{document}
